% GENERATED FILE. Edit canonical lessons, then run npm run generate:books.
% canonical-source-hash: fnv1a64:104b8ffb7a8e8451
% canonical-lessons: BN-C310-shongkho, BN-C310-prarthona-kora, BN-C310-bhokto, BN-C310-debi, BN-C310-pobitro

\chapter{Conch Shells and Prayers}
\label{ch:bn-conch-shells-and-prayers}
\hlchaptermodality{310}

\begin{chapteropening}
\textbf{By the end of this chapter:} \emph{I can talk about prayers, devotees and holy things.}

Complete the last lesson of chapter 310: I can talk about prayers, devotees and holy things.
\end{chapteropening}

\section[\texorpdfstring{\bn{শঙ্খ} (shôngkhô)}{শঙ্খ (shôngkhô)}]{\bn{শঙ্খ} (shôngkhô) — a conch shell}

\label{lesson:BN-C310-shongkho}



\begin{quote}
Before the new one: say the Bengali for a new account book, then the Bengali for a custom.
\end{quote}

\subsection*{You'll want to know: \bn{শঙ্খ}}
\textbf{\bn{শঙ্খ}} — \emph{shôngkhô} — \textquotedblleft{}a conch shell\textquotedblright{}.

\begin{quote}
\bn{সন্ধ্যায়} \bn{মন্দিরে} \bn{কেউ} \bn{শঙ্খ} \bn{বাজাল।}
\end{quote}

\begin{quote}
\emph{sôndhyāy mondire keu shôngkhô bājālô.}
\end{quote}

\begin{quote}
In the evening someone blew a conch at the temple.
\end{quote}

\begin{cousinweb}
\textbf{\bn{শঙ্খ}} (\emph{shôngkhô}) is Sanskrit \emph{śaṅkha}, a shell, above all the conch (Monier-Williams). Its \bn{ঙ্খ} is \emph{ngkh}; everyday speech wears the word down to \bn{শাঁখ}.
\end{cousinweb}

\begin{grammarlens}[title={What you've built}]
One more word for telling what happened.
\end{grammarlens}

\subsection*{Your turn}
\begin{itemize}
\raggedright
  \item \emph{Say it:} \emph{shôngkhô}
  \item \emph{Say it:} \emph{shôngkhô}, once more
  \item \emph{Say it:} \emph{sôndhyāy mondire keu shôngkhô bājālô.} — In the evening someone blew a conch at the temple.
  \item \emph{Recall:} say the Bengali for a new account book, then the Bengali for a custom, then say \emph{shôngkhô} again
\end{itemize}

\begin{tcolorbox}[breakable,colback=teal!4,colframe=teal!35!black,title={Before you move on}]
What does \bn{শঙ্খ} mean? (\textquotedblleft{}A conch shell\textquotedblright{}.) Say it once more.
\end{tcolorbox}
\section[\texorpdfstring{\bn{প্রার্থনা} \bn{করা} (prārthônā kôrā)}{প্রার্থনা করা (prārthônā kôrā)}]{\bn{প্রার্থনা} \bn{করা} (prārthônā kôrā) — to pray}

\label{lesson:BN-C310-prarthona-kora}



\begin{quote}
Before the new one: say the Bengali for a custom, then the Bengali for a conch shell.
\end{quote}

\subsection*{You'll want to know: \bn{প্রার্থনা} \bn{করা}}
\textbf{\bn{প্রার্থনা} \bn{করা}} — \emph{prārthônā kôrā} — \textquotedblleft{}to pray\textquotedblright{}.

\begin{quote}
\bn{বন্যার} \bn{সময়} \bn{সবাই} \bn{মন্দিরে} \bn{গিয়ে} \bn{প্রার্থনা} \bn{করল।}
\end{quote}

\begin{quote}
\emph{bônyār sômôy sôbāi mondire giye prārthônā kôrlô.}
\end{quote}

\begin{quote}
During the flood everyone went to the temple and prayed.
\end{quote}

Past (he/she …): \bn{তিনি} \bn{মন্দিরে} \bn{প্রার্থনা} \bn{করলেন}, \emph{tini mondire prārthônā kôrlen}.

\begin{cousinweb}
\textbf{\bn{প্রার্থনা} \bn{করা}} (\emph{prārthônā kôrā}) puts Sanskrit \emph{prārthanā}, 'a request, a petition; a prayer' (Monier-Williams), before \bn{করা}.
\end{cousinweb}

\begin{grammarlens}[title={What you've built}]
One more word for telling what happened.
\end{grammarlens}

\subsection*{Your turn}
\begin{itemize}
\raggedright
  \item \emph{Say it:} \emph{prārthônā kôrā}
  \item \emph{Say it:} \emph{prārthônā kôrā}, once more
  \item \emph{Say it:} \emph{bônyār sômôy sôbāi mondire giye prārthônā kôrlô.} — During the flood everyone went to the temple and prayed.
  \item \emph{Recall:} say the Bengali for a custom, then the Bengali for a conch shell, then say \emph{prārthônā kôrā} again
\end{itemize}

\begin{tcolorbox}[breakable,colback=teal!4,colframe=teal!35!black,title={Before you move on}]
What does \bn{প্রার্থনা} \bn{করা} mean? (\textquotedblleft{}To pray\textquotedblright{}.) Say it once more.
\end{tcolorbox}
\section[\texorpdfstring{\bn{ভক্ত} (bhôktô)}{ভক্ত (bhôktô)}]{\bn{ভক্ত} (bhôktô) — a devotee}

\label{lesson:BN-C310-bhokto}



\begin{quote}
Before the new one: say the Bengali for a conch shell, then the Bengali for to pray.
\end{quote}

\subsection*{You'll want to know: \bn{ভক্ত}}
\textbf{\bn{ভক্ত}} — \emph{bhôktô} — \textquotedblleft{}a devotee\textquotedblright{}.

\begin{quote}
\bn{পুজোর} \bn{দিন} \bn{মন্দিরে} \bn{অনেক} \bn{ভক্ত} \bn{এলেন।}
\end{quote}

\begin{quote}
\emph{pujor din mondire ônek bhôktô elen.}
\end{quote}

\begin{quote}
Many devotees came to the temple on the day of the puja.
\end{quote}

\begin{cousinweb}
\textbf{\bn{ভক্ত}} (\emph{bhôktô}) is Sanskrit \emph{bhakta}, first 'allotted, shared' and then 'devoted to, a worshipper' (Monier-Williams); \textbf{\bn{ভাত}}, cooked rice, comes down from the same word (CDIAL 9331).
\end{cousinweb}

\begin{grammarlens}[title={What you've built}]
One more word for telling what happened.
\end{grammarlens}

\subsection*{Your turn}
\begin{itemize}
\raggedright
  \item \emph{Say it:} \emph{bhôktô}
  \item \emph{Say it:} \emph{bhôktô}, once more
  \item \emph{Say it:} \emph{pujor din mondire ônek bhôktô elen.} — Many devotees came to the temple on the day of the puja.
  \item \emph{Recall:} say the Bengali for a conch shell, then the Bengali for to pray, then say \emph{bhôktô} again
\end{itemize}

\begin{tcolorbox}[breakable,colback=teal!4,colframe=teal!35!black,title={Before you move on}]
What does \bn{ভক্ত} mean? (\textquotedblleft{}A devotee\textquotedblright{}.) Say it once more.
\end{tcolorbox}
\section[\texorpdfstring{\bn{দেবী} (debī)}{দেবী (debī)}]{\bn{দেবী} (debī) — a goddess}

\label{lesson:BN-C310-debi}



\begin{quote}
Before the new one: say the Bengali for to pray, then the Bengali for a devotee.
\end{quote}

\subsection*{You'll want to know: \bn{দেবী}}
\textbf{\bn{দেবী}} — \emph{debī} — \textquotedblleft{}a goddess\textquotedblright{}.

\begin{quote}
\bn{গ্রামের} \bn{মানুষ} \bn{নদীর} \bn{দেবীকে} \bn{ভয়} \bn{পায়।}
\end{quote}

\begin{quote}
\emph{grāmer mānush nôdīr debīke bhôy pāy.}
\end{quote}

\begin{quote}
The villagers are afraid of the river goddess.
\end{quote}

\begin{cousinweb}
\textbf{\bn{দেবী}} (\emph{debī}) is Sanskrit \emph{devī}, 'a female deity, a goddess' (Monier-Williams), the feminine of \emph{deva}, a god.
\end{cousinweb}

\begin{grammarlens}[title={What you've built}]
One more word for telling what happened.
\end{grammarlens}

\subsection*{Your turn}
\begin{itemize}
\raggedright
  \item \emph{Say it:} \emph{debī}
  \item \emph{Say it:} \emph{debī}, once more
  \item \emph{Say it:} \emph{grāmer mānush nôdīr debīke bhôy pāy.} — The villagers are afraid of the river goddess.
  \item \emph{Recall:} say the Bengali for to pray, then the Bengali for a devotee, then say \emph{debī} again
\end{itemize}

\begin{tcolorbox}[breakable,colback=teal!4,colframe=teal!35!black,title={Before you move on}]
What does \bn{দেবী} mean? (\textquotedblleft{}A goddess\textquotedblright{}.) Say it once more.
\end{tcolorbox}
\section[\texorpdfstring{\bn{পবিত্র} (pôbitrô)}{পবিত্র (pôbitrô)}]{\bn{পবিত্র} (pôbitrô) — holy; sacred}

\label{lesson:BN-C310-pobitro}



\begin{quote}
Before the new one: say the Bengali for a devotee, then the Bengali for a goddess.
\end{quote}

\subsection*{You'll want to know: \bn{পবিত্র}}
\textbf{\bn{পবিত্র}} — \emph{pôbitrô} — \textquotedblleft{}holy; sacred\textquotedblright{}.

\begin{quote}
\bn{এই} \bn{পুকুরের} \bn{জল} \bn{পবিত্র।}
\end{quote}

\begin{quote}
\emph{ei pukurer jôl pôbitrô.}
\end{quote}

\begin{quote}
The water of this pond is holy.
\end{quote}

\begin{cousinweb}
\textbf{\bn{পবিত্র}} (\emph{pôbitrô}) is Sanskrit \emph{pavitra}, a means of purifying, a strainer, and so 'pure, holy' (Monier-Williams).
\end{cousinweb}

\begin{grammarlens}[title={What you've built}]
One more word for telling what happened.
\end{grammarlens}

\subsection*{Your turn}
\begin{itemize}
\raggedright
  \item \emph{Say it:} \emph{pôbitrô}
  \item \emph{Say it:} \emph{pôbitrô}, once more
  \item \emph{Say it:} \emph{ei pukurer jôl pôbitrô.} — The water of this pond is holy.
  \item \emph{Recall:} say the Bengali for a devotee, then the Bengali for a goddess, then say \emph{pôbitrô} again
\end{itemize}

\begin{tcolorbox}[breakable,colback=teal!4,colframe=teal!35!black,title={Before you move on}]
What does \bn{পবিত্র} mean? (\textquotedblleft{}Holy; sacred\textquotedblright{}.) Say it once more.
\end{tcolorbox}
